\documentclass[a4paper,12pt]{exam}
\usepackage[T1]{fontenc}
\usepackage{lmodern}
\usepackage{comment}
\usepackage{fancyvrb}
\usepackage{graphicx}
\usepackage{geometry}
\usepackage{longtable,booktabs}
\usepackage{graphicx}

%\setcounter{secnumdepth}{0}

\geometry{
  %includeheadfoot,
  margin=2.54cm
}

\title{Exercises --- Scala (Week Five)}
\author{\emph{Practice with classes}}
\date{Spring term 2017}

\begin{document}
\maketitle

\subsection*{Overview}

We now have enough of the syntax to have some fun playing with classes. 

\subsection*{Learning Objectives}

\begin{itemize}
	\item Practice with classes.
	\item Practice with case classes.
	\item Practice with companion classes.
\end{itemize}

\subsection*{Testing}

This exercise sheet does not come with any existing code or test harness --- 
you are required to write appropriate tests using the \verb!ScalaTest! framework.
You should follow \emph{good practice} with respect to \verb!sbt! folder structure
and build files (Hint: copy them from an earlier week).

\subsection*{The Questions}

\begin{questions}
\question
\begin{parts}
\part
Implement a \texttt{Counter} class. The constructor should take an \texttt{Int}. The methods \texttt{inc} and 
\texttt{dec} should increment and decrement the counter respectively returning a new \texttt{Counter}. 
Here’s an example of the usage:
\begin{Verbatim}
scala> new Counter(10).inc.dec.inc.inc.count
res02: Int = 12
\end{Verbatim}

\begin{solution}
\begin{Verbatim}
class Counter(val count: Int) { 
	def dec = new Counter(count - 1) 
	def inc = new Counter(count + 1)
}	
\end{Verbatim}	
Aside from practicing with classes and objects, this exercise has a second goal -—- to think about 
why \verb!inc!and \verb!dec! return a new \verb!Counter!, rather than updating 
the same counter directly.

Because \verb!val! fields are immutable, we need to come up with some other way of 
propagating the new value of \verb!count!. 
Methods that return new \verb!Counter! objects give us a way of returning new state 
without the side-effects of assignment. 
They also permit \emph{method chaining}, allowing us to write whole sequences of 
updates in a single expression

\textbf{Performance tip}: The use-case 
\begin{Verbatim}
new Counter(10).inc.dec.inc.inc.count 
\end{Verbatim}
actually creates five instances of \verb!Counter! before returning its final \verb!Int! value. 
You may be concerned about the extra memory and CPU overhead for such a simple 
calculation, but don’t be. Modern execution environments like the JVM render the 
extra overhead of this style of programming negligible in all but the most performance 
critical code.
\end{solution}


\part
Augment the \texttt{Counter} to 
allow the user can optionally pass an \texttt{Int} 
parameter to inc and dec. If the parameter is omitted it should default to 1.

\begin{solution}
The simplest solution is this:
\begin{Verbatim}
class Counter(val count: Int) {
	def dec(amount: Int = 1) = new Counter(count - amount) 
	def inc(amount: Int = 1) = new Counter(count + amount)
}	
\end{Verbatim}	
However, this adds parentheses to \verb!inc! and \verb!dec!. 
If we omit the parameter we now have to provide an empty pair of parentheses. 
We can work around this using method overloading to recreate our original parenthesis-
free methods. Note that overloading methods requires us to specify the return types:
\begin{Verbatim}
class Counter(val count: Int) { 
	def dec: Counter = dec()
	def inc: Counter = inc()
	def dec(amount: Int = 1): Counter = new Counter(count - amount)
	def inc(amount: Int = 1): Counter = new Counter(count + amount) 
}	
\end{Verbatim}
\end{solution}

\part
Reimplement \texttt{Counter} as a case class, using \texttt{copy} where appropriate. 
Additionally initialise \texttt{count} to a default value of \texttt{0}.
\begin{solution}
\begin{Verbatim}
case class Counter(count: Int = 0) {
	def dec = copy(count = count - 1)
	def inc = copy(count = count + 1)
	def adjust(adder: Adder) = copy(count = adder(count))
}	
\end{Verbatim}	
\end{solution}

\part
Here is a simple class called \texttt{Adder}:
\begin{Verbatim}
class Adder(amount: Int) {
  def add(in: Int) = in + amount
}
\end{Verbatim}
Extend \texttt{Counter} to add a method called \texttt{adjust}. 
This method should accept an \texttt{Adder} and return a new \texttt{Counter} with the result of 
applying the \texttt{Adder} to the \texttt{count}.
\begin{solution}
\begin{Verbatim}
class Counter(val count: Int) { 
	def dec = new Counter(count - 1) 
	def inc = new Counter(count + 1) 
	def adjust(adder: Adder) =
		new Counter(adder.add(count))
}	
\end{Verbatim}	
\end{solution}
\end{parts}

\question
\begin{parts}
\part
Implement a companion object for 
a \texttt{Person} class containing an \texttt{apply} method that accepts a 
whole name as a single string rather than individual first and last names.

Tip: you can split a \texttt{String} into an \texttt{Array} of components as follows:
\begin{Verbatim}
scala> val parts = "John Doe".split(" ")
parts: Array[String] = Array(John, Doe)

scala> parts(0)
res36: String = John
\end{Verbatim}
\begin{solution}
\begin{Verbatim}
object Person {
	def apply(name: String) = {
		val parts = name.split(" ")
		new Person(parts(0), parts(1)) 
	}
}	
\end{Verbatim}	
\end{solution}

\part
What happens when we define a companion object for a case class?

Take our \texttt{Person} class and turn it into a case class. 
Make sure you still have the companion object with the alternate \texttt{apply} method as well.
\begin{solution}
\begin{Verbatim}
case class Person(
			firstName: String, 
			lastName: String) { 
			
			def name = firstName + " " + lastName
}
			
object Person {
	def apply(name: String): Person = {
		val parts = name.split(" ")
	    apply(parts(0), parts(1))
	}
}	
\end{Verbatim}	
\end{solution}
\end{parts}

\question
\begin{parts}
\part
Write two classes, \texttt{Director} and \texttt{Film}, with fields and methods as follows:
\begin{itemize}
\item    \texttt{Director} should contain:
\begin{itemize}
\item        a field \texttt{firstName} of type \texttt{String}
\item        a field \texttt{lastName} of type \texttt{String}
\item        a field \texttt{yearOfBirth} of type \texttt{Int}
\item        a method called \texttt{name} that accepts no parameters and returns the full name
\end{itemize}
\item
    \texttt{Film} should contain:
    \begin{itemize}
    \item        a field \texttt{name} of type \texttt{String}
    \item    a field \texttt{yearOfRelease} of type \texttt{Int}
     \item   a field \texttt{imdbRating} of type \texttt{Double}
     \item   a field \texttt{director} of type \texttt{Director}
     \item   a method \texttt{directorsAge} that returns the age of the director at the time of release
      \item  a method \texttt{isDirectedBy} that accepts a \texttt{Director} as a parameter and 
      returns a \texttt{Boolean}
\end{itemize}
\end{itemize}
You will find appropriate demo data on the repo under the folder \verb!scala-exercises!;
you will need to adjust your constructors so that the code works without modification.

\begin{solution}
see next question\ldots	
\end{solution}

Implement a method of \texttt{Film} called \texttt{copy}. 
This method should accept the same parameters as the
constructor and create a new copy of the film. Give each parameter a default value so you can 
copy a film changing any subset of its values:
\begin{Verbatim}
highPlainsDrifter.copy(name = "L'homme des hautes plaines")
// returns Film("L'homme des hautes plaines", 1973, 7.7, /* etc */)

thomasCrownAffair.copy(yearOfRelease = 1968,
  director = new Director("Norman", "Jewison", 1926))
// returns Film("The Thomas Crown Affair", 1926, /* etc */)

inception.copy().copy().copy()
// returns a new copy of `inception`
\end{Verbatim}
\begin{solution}
see next question\ldots
\end{solution}

\part
Write companion objects for \texttt{Director} and \texttt{Film} as follows:
\begin{itemize}
\item
The \texttt{Director} companion object should contain:
\begin{itemize}
\item
    an \texttt{apply} method that accepts the same parameters as the constructor of the class 
    and returns a new \texttt{Director};
\item
    a method \texttt{older} that accepts two \texttt{Director}s and returns the oldest of the two.
\end{itemize}
\item
The \texttt{Film} companion object should contain:
\begin{itemize}
\item
    an \texttt{apply} method that accepts the same parameters as the constructor of the class and 
    returns a new \texttt{Film};
\item
    a method \texttt{highestRating} that accepts two \texttt{Film}s and returns the highest \texttt{imdbRating} 
    of the two;
\item
    a method \texttt{oldestDirectorAtTheTime} that accepts two \texttt{Film}s and returns the 
    \texttt{Director} who was oldest at the respective time of filming.

\end{itemize}
\end{itemize}

\begin{solution}
\begin{Verbatim}
class Director(
	val firstName: String, 
	val lastName: String, 
	val yearOfBirth: Int) {
	
	def name: String = s"$firstName $lastName"
	
	def copy(
		firstName: String = this.firstName,
		lastName: String = this.lastName,
		yearOfBirth: Int = this.yearOfBirth) =
			new Director(firstName, lastName, yearOfBirth)	
}

object Director {
	def apply(firstName: String, lastName: String, yearOfBirth: Int) =
		new Director(firstName, lastName, yearOfBirth)

	def older(director1: Director, director2: Director): Director =
		if (director1.yearOfBirth < director2.yearOfBirth) director1 else director2
}

class Film(
	val name: String,
	val yearOfRelease: Int, 
	val imdbRating: Double, 
	val director: Director) {
	
	def directorsAge =
		yearOfRelease - director.yearOfBirth
	
	def isDirectedBy(director: Director) = this.director == director
	
	def copy(
		name: String = this.name,
		yearOfRelease: Int = this.yearOfRelease, 
		imdbRating: Double = this.imdbRating,
		director: Director = this.director) =
			new Film(name, yearOfRelease, imdbRating, director)
}
		
object Film { 
	def apply(
		name: String,
		yearOfRelease: Int,
		imdbRating: Double,
		director: Director) =
			new Film(name, yearOfRelease, imdbRating, director)
	
	def newer(film1: Film, film2: Film): Film =
		if (film1.yearOfRelease < film2.yearOfRelease) film1 else film2
	
	def highestRating(film1: Film, film2: Film): Double = { 
		val rating1 = film1.imdbRating
		val rating2 = film2.imdbRating
		if (rating1 > rating2) rating1 else rating2
	}
		
	def oldestDirectorAtTheTime(film1: Film, film2: Film): Director =
		if (film1.directorsAge > film2.directorsAge) film1.director else film2.director
}
\end{Verbatim}	
\end{solution}

\part
We can dispose of much of the \textit{boilerplate} by converting the \texttt{Director} and \texttt{Film} classes
to \textit{case classes}. Do this conversion and work out what code we can remove.

\end{parts}
\end{questions}
\end{document}

\question
The similarity in naming of classes and companion objects tends to cause confusion for new Scala developers. 
When reading a block of code it is important to know which parts refer to a class or type and which parts refer 
to a singleton object or value.

This is the inspiration for the new hit quiz, \textit{Type or Value?}, which we will be piloting below. 
In each case identify whether the word \texttt{Film} refers to the type or value:
\begin{parts}
\part
\begin{Verbatim}
val prestige: Film = bestFilmByChristopherNolan()
\end{Verbatim}
\begin{solution}
Type
\end{solution}

\part
\begin{Verbatim}
new Film("Last Action Hero", 1993, mcTiernan)
\end{Verbatim}
\begin{solution}
Type
\end{solution}


\part
\begin{Verbatim}
Film("Last Action Hero", 1993, mcTiernan)
\end{Verbatim}
\begin{solution}
Value
\end{solution}

\part
\begin{Verbatim}
Film.newer(highPlainsDrifter, thomasCrownAffair)
\end{Verbatim}
\begin{solution}
Value
\end{solution}

\part
\begin{Verbatim}
Film.type
\end{Verbatim}
\begin{solution}
Value
\end{solution}
\end{parts}